% Options for packages loaded elsewhere
\PassOptionsToPackage{unicode}{hyperref}
\PassOptionsToPackage{hyphens}{url}
\PassOptionsToPackage{dvipsnames,svgnames,x11names}{xcolor}
%
\documentclass[
  11pt,
]{article}
\usepackage{amsmath,amssymb}
\usepackage{iftex}
\ifPDFTeX
  \usepackage[T1]{fontenc}
  \usepackage[utf8]{inputenc}
  \usepackage{textcomp} % provide euro and other symbols
\else % if luatex or xetex
  \usepackage{unicode-math} % this also loads fontspec
  \defaultfontfeatures{Scale=MatchLowercase}
  \defaultfontfeatures[\rmfamily]{Ligatures=TeX,Scale=1}
\fi
\usepackage{lmodern}
\ifPDFTeX\else
  % xetex/luatex font selection
\fi
% Use upquote if available, for straight quotes in verbatim environments
\IfFileExists{upquote.sty}{\usepackage{upquote}}{}
\IfFileExists{microtype.sty}{% use microtype if available
  \usepackage[]{microtype}
  \UseMicrotypeSet[protrusion]{basicmath} % disable protrusion for tt fonts
}{}
\makeatletter
\@ifundefined{KOMAClassName}{% if non-KOMA class
  \IfFileExists{parskip.sty}{%
    \usepackage{parskip}
  }{% else
    \setlength{\parindent}{0pt}
    \setlength{\parskip}{6pt plus 2pt minus 1pt}}
}{% if KOMA class
  \KOMAoptions{parskip=half}}
\makeatother
\usepackage{xcolor}
\usepackage[margin=1in]{geometry}
\usepackage{longtable,booktabs,array}
\usepackage{calc} % for calculating minipage widths
% Correct order of tables after \paragraph or \subparagraph
\usepackage{etoolbox}
\makeatletter
\patchcmd\longtable{\par}{\if@noskipsec\mbox{}\fi\par}{}{}
\makeatother
% Allow footnotes in longtable head/foot
\IfFileExists{footnotehyper.sty}{\usepackage{footnotehyper}}{\usepackage{footnote}}
\makesavenoteenv{longtable}
\setlength{\emergencystretch}{3em} % prevent overfull lines
\providecommand{\tightlist}{%
  \setlength{\itemsep}{0pt}\setlength{\parskip}{0pt}}
\setcounter{secnumdepth}{-\maxdimen} % remove section numbering
\ifLuaTeX
  \usepackage{selnolig}  % disable illegal ligatures
\fi
\IfFileExists{bookmark.sty}{\usepackage{bookmark}}{\usepackage{hyperref}}
\IfFileExists{xurl.sty}{\usepackage{xurl}}{} % add URL line breaks if available
\urlstyle{same}
\hypersetup{
  pdfauthor={Stefano Alimonti --- stefalimonti@gmail.com},
  colorlinks=true,
  linkcolor={blue},
  filecolor={Maroon},
  citecolor={Blue},
  urlcolor={blue},
  pdfcreator={LaTeX via pandoc}}

\author{Stefano Alimonti --- stefalimonti@gmail.com}
\date{March 2026}

\begin{document}

{
\hypersetup{linkcolor=black}
\setcounter{tocdepth}{2}
\tableofcontents
}
\section{Eigenvalue Statistics of Carry Companion Matrices: A Markov-Driven
GOE↔GUE Transition in Sparse Non-Hermitian
Ensembles}\label{eigenvalue-statistics-of-carry-companion-matrices-a-markov-driven-goegue-transition-in-sparse-non-hermitian-ensembles}

\textbf{Author:} Stefano Alimonti \textbf{Affiliation:} Independent Researcher
\textbf{Date:} March 2026

\begin{center}\rule{0.5\linewidth}{0.5pt}\end{center}

\newpage

\subsection{Abstract}\label{abstract}

We study eigenvalue spacing statistics for companion matrices arising from carry
polynomials in positional multiplication. Individual carry matrices exhibit
\textbf{GOE-like} nearest-neighbor spacing (orthogonal symmetry class). We
identify the mechanism: Markov correlations between adjacent entries of the last
column create an effective rank reduction in the perturbation space, quantified
by a one-parameter interpolation that maps the full GOE↔GUE transition.
Replacing carries with i.i.d. entries of the same marginal distribution shifts
statistics to \textbf{GUE} (unitary class).

Three analytical results underpin the transition: (1) the carry chain possesses
a provably positive lag-1 autocorrelation (Proposition 1); (2) a stationary
binary Markov chain with correlation \(\rho\) produces a covariance matrix whose
normalized effective rank is

\[g(\rho) = \frac{1-\rho^2}{1+\rho^2},\]

yielding \(g(1/2) = 3/5\) at the Diaconis--Fulman spectral gap value
(Proposition 2); and (3) the real eigenvalue fraction increases monotonically
with \(\rho\), anti-correlated with the effective rank (Observation 3).

Four independent statistics confirm the transition --- the spacing ratio
\(\langle\tilde{r}\rangle\), the number variance \(\Sigma^2(L)\), the
eigenvector orthogonality structure, and the real eigenvalue fraction --- with a
computational bound

\[\beta_{\mathrm{eff}} \leq 2 - \varepsilon(D), \qquad \varepsilon > 0,\]

at 99.9\% bootstrap confidence.

\textbf{Keywords:} random matrix theory, GOE-GUE transition, carry matrices,
eigenvalue statistics, Markov correlation, companion matrices, effective rank

\textbf{MSC 2020:} 15B52, 60B20, 11A63

\begin{center}\rule{0.5\linewidth}{0.5pt}\end{center}

\newpage

\subsection{1. Introduction}\label{introduction}

\subsubsection{1.1 Motivation}\label{motivation}

The spectral statistics of random matrices provide a bridge between arithmetic
structure and physical universality classes. Montgomery {[}1{]} proved that
Riemann zeta zeros exhibit GUE pair correlation; the Katz--Sarnak philosophy
{[}2{]} extends this correspondence to families of \(L\)-functions. In this
paper, we study whether the carry companion matrices --- sparse non-Hermitian
matrices arising from positional multiplication --- exhibit random matrix
statistics, and if so, which universality class.

\subsubsection{1.2 RMT Context and
Universality}\label{rmt-context-and-universality}

The spectral statistics of sparse non-Hermitian random matrices --- specifically
companion matrices where only one column contains random entries --- remain
challenging for rigorous RMT. While bulk universality has been established for
broad classes of Wigner matrices (Erdős, Schlein, and Yau {[}3{]}; Erdős and Yau
{[}4{]}), the companion structure falls outside these standard frameworks. Tao
and Vu {[}5{]} demonstrated that certain short-range correlations do not break
local universality for non-Hermitian matrices. Our results present a contrasting
phenomenon for sparse companion matrices: structured Markovian correlations
\emph{do} alter the universality class, driving a smooth continuous transition
--- analogous to the interpolating \(\beta\)-ensembles of Dumitriu and Edelman
{[}6{]} --- from a GUE-like uncorrelated state to a GOE-like strongly correlated
state. This positions the carry sequence as a naturally occurring operator
driving a rank-reducing RMT transition.

\subsubsection{1.3 Summary of Results}\label{summary-of-results}

We prove three analytical propositions and establish three computational
results:

\begin{enumerate}
\def\labelenumi{\arabic{enumi}.}
\tightlist
\item
  \textbf{Proposition 1} (§2.1): The carry chain has provably positive lag-1
  autocorrelation, with
\end{enumerate}

\[\mathrm{Corr}(c_2, c_3) = 1/\sqrt{2}\]

exactly, establishing the minimum bulk correlation. 2. \textbf{Proposition 2}
(§2.2): The normalized effective rank of the Markov covariance matrix is
\(g(\rho) = (1-\rho^2)/(1+\rho^2)\), strictly decreasing, with \(g(1/2) = 3/5\)
at the Diaconis--Fulman spectral gap value. 3. \textbf{Observation 3} (§4.5):
The real eigenvalue fraction \(f_{\mathrm{real}}(\rho)\) is monotonically
increasing in \(\rho\), anti-correlated with \(g(\rho)\) (\(r = -0.93\)),
providing the most direct TRS signature. 4. \textbf{Computational bound} (§4.3):

\[\beta_{\mathrm{eff}} \leq 2 - \varepsilon(D)\]

with \(\varepsilon > 0\) at 99.9\% bootstrap confidence for all tested \(D\). 5.
\textbf{Monotone GOE↔GUE interpolation} (§4.1):
\(\beta_{\mathrm{eff}}(\lambda)\) varies continuously from \({\approx}1.8\)
(near-GUE) to \({\approx}1.0\) (GOE). 6. \textbf{Four-way confirmation} (§3,
§4): Spacing ratio, number variance, eigenvector structure, and real eigenvalue
fraction independently confirm the transition.

\subsubsection{Important Caveats}\label{important-caveats}

All computational results are \textbf{empirical observations} established with
high statistical significance but without rigorous mathematical proof. Carry
matrices are sparse companion matrices with integer entries, not members of any
standard random matrix ensemble. The parallel with Riemann zero statistics is
suggestive but not formal.

\begin{center}\rule{0.5\linewidth}{0.5pt}\end{center}

\newpage

\subsection{2. Analytical Foundations}\label{analytical-foundations}

\subsubsection{2.1 Carry Coefficient
Autocorrelation}\label{carry-coefficient-autocorrelation}

The GOE↔GUE transition is driven by the internal correlations of the carry
chain. Using the exact second-moment structure developed in {[}F{]}, we
establish:

\textbf{Proposition 1} (Positive carry autocorrelation). For base-2
multiplication with fixed MSBs \(g_0 = h_0 = 1\) and remaining bits i.i.d.
\(\mathrm{Ber}(1/2)\):

\emph{(a) The lag-1 autocorrelation at the first non-trivial position is}

\[\mathrm{Corr}(c_2, c_3) = \frac{1}{\sqrt{2}} \approx 0.7071.\]

\begin{enumerate}
\def\labelenumi{(\alph{enumi})}
\setcounter{enumi}{1}
\item
  For all \(j \geq 2\) in the bulk, \(\mathrm{Cov}(c_j, c_{j+1}) > 0\).
\item
  The bulk autocorrelation is monotonically increasing with \(j\).
\end{enumerate}

\textbf{Proof.} (a) From the carry recurrence \(c_1 = 0\) and
\(\mathrm{conv}_1 = g_1 + h_1\), we have
\(c_2 = \lfloor(g_1 + h_1)/2\rfloor \in \{0,1\}\) with \(P(c_2 = 0) = 3/4\),
\(P(c_2 = 1) = 1/4\). Hence \(E[c_2] = 1/4\), \(\mathrm{Var}(c_2) = 3/16\) {[}F,
Theorem 3{]}.

For \(c_3\): when \(c_2 = 1\) (i.e., \(g_1 = h_1 = 1\)), we have

\[\mathrm{conv}_2 = h_2 + g_1 h_1 + g_2 = h_2 + 1 + g_2\]

and

\[\mathrm{total}_2 = \mathrm{conv}_2 + 1,\]

yielding \(E[c_3 \mid c_2 = 1] = 5/4\). By the law of total expectation and the
exact moment \(E[c_3] = 1/2\) {[}F, Theorem 3{]}, \(\mathrm{Var}(c_3) = 3/8\):

\[\mathrm{Cov}(c_2, c_3) = E[c_2 c_3] - E[c_2]E[c_3] = \frac{5}{16} - \frac{1}{4}\cdot\frac{1}{2} = \frac{3}{16}.\]

Therefore

\[\mathrm{Corr}(c_2, c_3) = \frac{3/16}{\sqrt{3/16 \cdot 3/8}} = \frac{3/16}{3/(8\sqrt{2})} = \frac{1}{\sqrt{2}}.\]

\(\square\)

\begin{enumerate}
\def\labelenumi{(\alph{enumi})}
\setcounter{enumi}{1}
\tightlist
\item
  The carry recurrence
\end{enumerate}

\[c_{j+1} = \lfloor(\mathrm{conv}_j + c_j)/2\rfloor\]

is weakly increasing in \(c_j\) for fixed \(\mathrm{conv}_j\) (increasing
\(c_j\) by 1 either preserves or increases \(c_{j+1}\) by 1). Hence \(c_{j+1}\)
is stochastically increasing in \(c_j\), giving
\(\mathrm{Cov}(c_j, c_{j+1}) \geq 0\). Strict positivity follows because, with
probability \(1/2\) (by the Parity Lemma {[}F, Lemma B{]}), \(\mathrm{total}_j\)
is odd, in which case the unit increase in \(c_j\) flips \(c_{j+1}\) from
\(\lfloor t/2 \rfloor\) to \(\lfloor(t+1)/2\rfloor = \lfloor t/2\rfloor + 1\).
\(\square\)

\textbf{Observation} (monotonicity; numerical). The bulk autocorrelation
\(\mathrm{Corr}(c_j, c_{j+1})\) is monotonically increasing with \(j\) in all
tested cases: verified through exact rational enumeration up to \(K = 8\) (14
free bits, \(2^{14}\) configurations) and by Monte Carlo to \(K = 64\). A full
analytical proof of monotonicity remains open.

\textbf{Remark} (Three types of correlation). It is important to distinguish: -
(i) The raw carry autocorrelation \(\mathrm{Corr}(c_j, c_{j+1})\), which starts
at \(1/\sqrt{2}\) and increases toward 1 as \(j \to \infty\) (because both means
grow linearly); - (ii) The companion matrix coefficient correlation
\(\mathrm{Corr}(q_i, q_{i+1})\) for the quotient polynomial \(Q = C/(x-2)\),
which inherits the carry structure; - (iii) The Diaconis--Fulman spectral gap
\(\gamma = 1 - 1/b = 1/2\) for base 2 {[}A; 7{]}, which governs the \emph{rate
of decorrelation} in the carry Markov chain.

The binary model of §4.4 uses \(\rho = 1/2\) as its lag-1 correlation parameter,
matching the spectral gap (iii), not the raw correlation (i). This is the
correct physical matching: the spectral gap controls the effective
dimensionality of the perturbation vector (Proposition 2).

\subsubsection{2.2 Effective Rank of the Markov Correlation
Model}\label{effective-rank-of-the-markov-correlation-model}

To isolate the effect of correlation \(\rho\) on spectral statistics, we study
the covariance matrix of a stationary binary Markov chain. Let \(\Sigma(\rho)\)
be the \(D \times D\) covariance matrix of a length-\(D\) stationary binary
Markov chain with \(\mathrm{Ber}(1/2)\) marginals and lag-1 correlation
\(\rho\).

\textbf{Proposition 2} (Effective rank formula). The covariance matrix
\(\Sigma(\rho)\) is a Kac--Murdock--Szegő Toeplitz matrix with
\(\Sigma_{ij} = \sigma^2 \rho^{|i-j|}\) (\(\sigma^2 = 1/4\)). Its normalized
effective rank (participation ratio)

\[r_{\mathrm{eff}} = \frac{(\mathrm{Tr}\,\Sigma)^2}{D \cdot \mathrm{Tr}\,\Sigma^2}\]

\emph{satisfies:}

\emph{(a)} \(r_{\mathrm{eff}}(0) = 1\) \emph{(uncorrelated limit, full rank);}

\emph{(b)}

\[g(\rho) := \lim_{D \to \infty} r_{\mathrm{eff}} = \frac{1 - \rho^2}{1 + \rho^2},\]

which is strictly decreasing on \([0, 1)\) with
\(g'(\rho) = -4\rho/(1+\rho^2)^2 < 0\) for \(\rho > 0\);

\emph{(c)} \(g(1/2) = 3/5 = 0.6\) \emph{exactly. At the carry chain spectral gap
value, the perturbation vector retains only 60\% of its nominal degrees of
freedom.}

\textbf{Proof.} The eigenvalues of the KMS matrix are given by the
Grenander--Szegő formula {[}8{]}:

\[\lambda_k = \sigma^2 \frac{1-\rho^2}{1 - 2\rho\cos\!\bigl(\frac{k\pi}{D+1}\bigr) + \rho^2}, \qquad k = 1, \ldots, D.\]

In the \(D \to \infty\) limit, the Szegő distribution theorem converts the
discrete sums to integrals over the spectral density
\(f(\theta) = \sigma^2 P_\rho(\theta)\), where
\(P_\rho(\theta) = (1-\rho^2)/(1-2\rho\cos\theta+\rho^2)\) is the Poisson
kernel. We compute:

\[\frac{1}{\pi}\int_0^\pi P_\rho(\theta)\,d\theta = 1, \qquad \frac{1}{\pi}\int_0^\pi P_\rho(\theta)^2\,d\theta = \frac{1+\rho^2}{1-\rho^2}.\]

The first identity is standard. The second follows from Parseval's theorem
applied to the Fourier expansion
\(P_\rho(\theta) = \sum_{n=-\infty}^{\infty} \rho^{|n|} e^{in\theta}\):

\[\frac{1}{2\pi}\int_0^{2\pi} P_\rho(\theta)^2\,d\theta = \sum_{n=-\infty}^{\infty} \rho^{2|n|} = 1 + \frac{2\rho^2}{1-\rho^2} = \frac{1+\rho^2}{1-\rho^2}.\]

The normalized effective rank in the continuous limit is then:

\[g(\rho) = \frac{\bigl[\int_0^\pi f(\theta)\,d\theta\bigr]^2}{\pi\int_0^\pi f(\theta)^2\,d\theta} = \frac{(\sigma^2\pi)^2}{\pi \cdot \sigma^4 \cdot \pi(1+\rho^2)/(1-\rho^2)} = \frac{1-\rho^2}{1+\rho^2}.\]

The strict decrease follows from \(g'(\rho) = -4\rho/(1+\rho^2)^2 < 0\) for
\(\rho > 0\). Evaluating: \(g(1/2) = (3/4)/(5/4) = 3/5\). The finite-\(D\)
correction is \(O(1/D)\), as verified numerically for \(D = 10\) to \(D = 5000\)
(C13, Prop. 2). \(\square\)

\textbf{Interpretation.} The effective rank reduction from 1 to \(3/5\) at
\(\rho = 1/2\) means the companion matrix \(M = J + v \cdot e_D^T\) (where \(J\)
is the shift matrix and \(v\) is the last column) has \(40\%\) fewer independent
perturbation directions. Fewer independent directions reduce the level repulsion
from quadratic (\(\beta = 2\), GUE) toward linear (\(\beta = 1\), GOE). This
rank reduction is the physical mechanism driving the universality class
transition.

\begin{center}\rule{0.5\linewidth}{0.5pt}\end{center}

\newpage

\subsection{3. The Key Result: Markov Correlation Imposes
GOE}\label{the-key-result-markov-correlation-imposes-goe}

\subsubsection{3.1 Controlled Comparison}\label{controlled-comparison}

Two companion matrix ensembles compared:

\begin{longtable}[]{@{}llll@{}}
\toprule\noalign{}
Ensemble & \(L^2\)(GOE) & \(L^2\)(GUE) & Best fit \\
\midrule\noalign{}
\endhead
\bottomrule\noalign{}
\endlastfoot
Markov-correlated (actual carries) & 0.084 & 0.109 & \textbf{GOE} \\
i.i.d. (same marginals, no correlation) & 0.146 & 0.119 & \textbf{GUE} \\
\end{longtable}

\subsubsection{3.2 Interpretation}\label{interpretation}

The empirical carry ensemble (with its actual Markov-induced dependence
structure) is GOE-like, while i.i.d. controls are GUE-like. In the simplified
binary Markov interpolation model of §4.4, increasing \(\rho\) moves the system
from sub-GOE through a GOE crossing near \(\rho \approx 1/2\) and then toward
GOE-GUE for larger \(\rho\). The operative mechanism is therefore not ``more
correlation always means more GOE'', but the specific structured dependence
regime realized by carry matrices, where the effective-rank reduction of
Proposition 2 and real-eigenvalue enhancement align near the GOE crossing.

\subsubsection{3.3 Level Spacing Ratio Test
(C07)}\label{level-spacing-ratio-test-c07}

The level spacing ratio \(\tilde{r}_n = \min(s_n, s_{n+1})/\max(s_n, s_{n+1})\)
(Atas et al.~{[}9{]}) provides a robust test that does not require spectral
unfolding. Known reference values:
\(\langle\tilde{r}\rangle_{\text{Poisson}} \approx 0.386\),
\(\langle\tilde{r}\rangle_{\text{GOE}} \approx 0.531\),
\(\langle\tilde{r}\rangle_{\text{GUE}} \approx 0.600\).

Markov vs i.i.d. comparison (20-bit primes, \(D \approx 38\)):

\begin{longtable}[]{@{}
  >{\raggedright\arraybackslash}p{(\columnwidth - 8\tabcolsep) * \real{0.1333}}
  >{\raggedright\arraybackslash}p{(\columnwidth - 8\tabcolsep) * \real{0.3600}}
  >{\raggedright\arraybackslash}p{(\columnwidth - 8\tabcolsep) * \real{0.1067}}
  >{\raggedright\arraybackslash}p{(\columnwidth - 8\tabcolsep) * \real{0.2800}}
  >{\raggedright\arraybackslash}p{(\columnwidth - 8\tabcolsep) * \real{0.1200}}@{}}
\toprule\noalign{}
\begin{minipage}[b]{\linewidth}\raggedright
Ensemble
\end{minipage} & \begin{minipage}[b]{\linewidth}\raggedright
\(\langle\tilde{r}\rangle\)
\end{minipage} & \begin{minipage}[b]{\linewidth}\raggedright
95\% CI
\end{minipage} & \begin{minipage}[b]{\linewidth}\raggedright
\(N_{\text{ratios}}\)
\end{minipage} & \begin{minipage}[b]{\linewidth}\raggedright
Closest
\end{minipage} \\
\midrule\noalign{}
\endhead
\bottomrule\noalign{}
\endlastfoot
Markov-correlated (actual) & 0.511 & {[}0.509, 0.513{]} & 66,402 &
\textbf{GOE} \\
i.i.d. (shuffled) & 0.586 & {[}0.584, 0.587{]} & 217,471 & \textbf{GUE} \\
\end{longtable}

The separation is \(0.075\) (\(Z = 58.6\), highly significant), exceeding the
GOE-GUE reference gap of \(0.069\).

\textbf{Dimension dependence (base 2):}

\begin{longtable}[]{@{}llll@{}}
\toprule\noalign{}
Bits & Mean \(D\) & \(\langle\tilde{r}\rangle\) & Closest \\
\midrule\noalign{}
\endhead
\bottomrule\noalign{}
\endlastfoot
10 & 18 & 0.579 & GUE \\
14 & 26 & 0.543 & GOE \\
18 & 34 & 0.526 & GOE \\
22 & 42 & 0.508 & GOE \\
26 & 50 & 0.495 & GOE \\
\end{longtable}

The spacing ratio converges toward GOE as \(D\) grows but at \(D = 50\) sits
slightly below the GOE value. This sub-GOE effect may reflect the sparse
companion structure deviating from standard \(\beta\)-ensembles (see §4.3
Caveats). Higher-dimensional tests (\(D > 100\)) would resolve this.

\subsubsection{3.4 Requirements for a Rigorous
Proof}\label{requirements-for-a-rigorous-proof}

A mathematical proof that Markov-correlated companion matrices exhibit GOE would
need:

\begin{enumerate}
\def\labelenumi{\arabic{enumi}.}
\tightlist
\item
  A universality theorem for sparse companion matrices with correlated entries
  (no such result exists in RMT)
\item
  A quantitative bound relating the effective rank reduction (Proposition 2) to
  the level repulsion exponent \(\beta\)
\item
  An explanation of the sub-GOE effect at large \(D\) within the
  \(\beta\)-ensemble framework
\end{enumerate}

\begin{center}\rule{0.5\linewidth}{0.5pt}\end{center}

\newpage

\subsection{4. Mechanism: Quantitative GOE↔GUE
Transition}\label{mechanism-quantitative-goegue-transition}

\subsubsection{4.1 Interpolation Experiment
(C08)}\label{interpolation-experiment-c08}

We construct a one-parameter family of companion matrices interpolating between
full Markov correlation (\(\lambda = 1\), actual carries) and i.i.d. entries
(\(\lambda = 0\), shuffled carries with same marginal). At intermediate
\(\lambda\), a fraction \(1 - \lambda\) of the carry positions are randomly
reshuffled.

\begin{longtable}[]{@{}
  >{\raggedright\arraybackslash}p{(\columnwidth - 8\tabcolsep) * \real{0.1392}}
  >{\raggedright\arraybackslash}p{(\columnwidth - 8\tabcolsep) * \real{0.3418}}
  >{\raggedright\arraybackslash}p{(\columnwidth - 8\tabcolsep) * \real{0.1013}}
  >{\raggedright\arraybackslash}p{(\columnwidth - 8\tabcolsep) * \real{0.3038}}
  >{\raggedright\arraybackslash}p{(\columnwidth - 8\tabcolsep) * \real{0.1139}}@{}}
\toprule\noalign{}
\begin{minipage}[b]{\linewidth}\raggedright
\(\lambda\)
\end{minipage} & \begin{minipage}[b]{\linewidth}\raggedright
\(\langle\tilde{r}\rangle\)
\end{minipage} & \begin{minipage}[b]{\linewidth}\raggedright
95\% CI
\end{minipage} & \begin{minipage}[b]{\linewidth}\raggedright
\(\beta_{\mathrm{eff}}\)
\end{minipage} & \begin{minipage}[b]{\linewidth}\raggedright
Closest
\end{minipage} \\
\midrule\noalign{}
\endhead
\bottomrule\noalign{}
\endlastfoot
0.00 & 0.5882 & {[}0.587, 0.589{]} & 1.84 & GUE \\
0.25 & 0.5716 & {[}0.571, 0.573{]} & 1.59 & GUE \\
0.50 & 0.5491 & {[}0.548, 0.550{]} & 1.27 & GOE \\
0.75 & 0.5316 & {[}0.531, 0.533{]} & 1.01 & GOE \\
1.00 & 0.5165 & {[}0.514, 0.519{]} & \(\leq 1\) & GOE \\
\end{longtable}

The transition is \textbf{smooth and monotone}:
\(\beta_{\mathrm{eff}}(\lambda)\) decreases continuously from \({\approx}1.84\)
at \(\lambda = 0\) to \({\approx}1.0\) at \(\lambda \approx 0.75\), then
saturates. At \(\lambda = 1\), \(\langle\tilde{r}\rangle\) is slightly
\emph{below} the GOE value (\(0.517\) vs \(0.531\)), a sub-GOE effect consistent
with the sparse companion structure.

\textbf{Note:} \(\beta_{\mathrm{eff}}\) does not reach \(2.0\) at
\(\lambda = 0\) because the companion matrix structure (sub-diagonal of 1s, all
randomness in one column) constrains the ensemble even when entries are i.i.d.

\subsubsection{4.2 Eigenvector Orthogonality Mechanism
(C09)}\label{eigenvector-orthogonality-mechanism-c09}

For a diagonalizable matrix \(M = V \Lambda V^{-1}\), the eigenvector structure
reveals the symmetry mechanism:

\begin{longtable}[]{@{}lllll@{}}
\toprule\noalign{}
Metric & Markov & i.i.d. & \(Z\)-score & Direction \\
\midrule\noalign{}
\endhead
\bottomrule\noalign{}
\endlastfoot
\(\kappa(V)\) (median) & 131 & 21 & 1.6 & Markov \emph{less} orthogonal \\
Gram off-diagonal & 0.78 & 0.28 & 121 & Markov \emph{less} orthogonal \\
Fraction real eigenvalues & 0.049 & 0.038 & 9.6 & Markov \emph{more} real \\
\end{longtable}

Markov-correlated entries make eigenvectors \emph{less} orthogonal (\(\kappa\)
is \(6\times\) larger), but simultaneously produce \emph{more real eigenvalues}
(\(+30\%\), \(Z = 9.6\)). The correlation
\(\rho(\log_{10} \kappa, \langle\tilde{r}\rangle) = -0.36\) is
\textbf{negative}: higher \(\kappa\) (less orthogonal) correlates with
\emph{lower} \(\langle\tilde{r}\rangle\) (more GOE-like).

\textbf{Interpretation.} The GOE mechanism is not ``Markov makes eigenvectors
orthogonal.'' Rather, the Markov carry recurrence imposes a \textbf{structured}
pattern of non-orthogonality --- constrained by the carry chain --- that
effectively preserves time-reversal symmetry. The \(\kappa(V)\) ratio
(Markov/i.i.d.) \textbf{grows with dimension}: \(3.6\times\) at \(D = 18\),
\(7.3\times\) at \(D = 50\), consistent with the spacing ratio converging deeper
into the GOE regime at larger \(D\).

\subsubsection{4.3 Computational Bound:}\label{computational-bound}

\[\beta_{\mathrm{eff}} \leq 2 - \varepsilon\]

(C11)

The Atas et al.~{[}9{]} Wigner surmise provides an analytic map
\(\langle\tilde{r}\rangle(\beta)\) valid for all \(\beta > 0\):

\[p_\beta(r) \propto \frac{(r + r^2)^\beta}{(1 + r + r^2)^{1 + 3\beta/2}}, \qquad \langle\tilde{r}\rangle(\beta) = \int_0^1 r\,\tilde{p}_\beta(r)\,dr.\]

Inverting this map on bootstrap confidence intervals yields a computational
bound.

\textbf{Proposition} (Computational \(\beta_{\mathrm{eff}}\) bound). *For
\(D\)-dimensional companion matrices with base-2 Markov-correlated carry
entries, the effective Dyson index satisfies

\[\beta_{\mathrm{eff}} \leq 2 - \varepsilon(D)\]

where \(\varepsilon(D) > 0\) at 99.9\% bootstrap confidence (\(B = 10{,}000\)
resamples):*

\begin{longtable}[]{@{}
  >{\raggedright\arraybackslash}p{(\columnwidth - 8\tabcolsep) * \real{0.1176}}
  >{\raggedright\arraybackslash}p{(\columnwidth - 8\tabcolsep) * \real{0.2353}}
  >{\raggedright\arraybackslash}p{(\columnwidth - 8\tabcolsep) * \real{0.3725}}
  >{\raggedright\arraybackslash}p{(\columnwidth - 8\tabcolsep) * \real{0.1569}}
  >{\raggedright\arraybackslash}p{(\columnwidth - 8\tabcolsep) * \real{0.1176}}@{}}
\toprule\noalign{}
\begin{minipage}[b]{\linewidth}\raggedright
\(\bar{D}\)
\end{minipage} & \begin{minipage}[b]{\linewidth}\raggedright
\(\beta_{\mathrm{eff}}\)
\end{minipage} & \begin{minipage}[b]{\linewidth}\raggedright
\(\beta_{\mathrm{upper}}\) (99.9\% CI)
\end{minipage} & \begin{minipage}[b]{\linewidth}\raggedright
\(\varepsilon\)
\end{minipage} & \begin{minipage}[b]{\linewidth}\raggedright
\(Z\) vs GUE
\end{minipage} \\
\midrule\noalign{}
\endhead
\bottomrule\noalign{}
\endlastfoot
18 & 1.56 & 1.64 & 0.37 & 17.7 \\
26 & 1.15 & 1.20 & 0.80 & 49.6 \\
34 & 0.85 & 0.89 & 1.12 & 83.0 \\
42 & 0.69 & 0.73 & 1.27 & 93.8 \\
50 & 0.60 & 0.63 & 1.37 & 105.7 \\
\end{longtable}

The bound \(\varepsilon(D)\) is \textbf{monotonically increasing} with \(D\):
the deviation from GUE strengthens with matrix size and is not a finite-size
artifact. The i.i.d. control ensemble yields
\(\beta_{\mathrm{eff}} \approx 1.7\)--\(2.2\) (near-GUE) across all \(D\),
isolating the Markov correlation as the sole mechanism.

\textbf{Caveats.} (1) The Wigner surmise is approximate; its
\(\langle\tilde{r}\rangle\) values deviate from exact ensemble averages by
\({\sim}1\%\) at \(\beta = 1, 2\). (2) At \(D \geq 34\),
\(\beta_{\mathrm{eff}} < 1\) (``sub-GOE''), which lies outside standard
\(\beta\)-ensemble theory and reflects the sparse companion structure rather
than a physical \(\beta < 1\) universality class. (3) This is a computational
result, not a mathematical proof.

\subsubsection{4.4 Mechanism: Effective Rank Reduction
(C12)}\label{mechanism-effective-rank-reduction-c12}

A simplified binary model isolates the mechanism analytically. Consider
\(D\)-dimensional companion matrices with last-column entries
\((a_0, \ldots, a_{D-2}) \in \{0,1\}\) from a symmetric Markov chain with lag-1
correlation \(\rho \in [0,1]\), and \(a_{D-1} = 1\) (ULC).

\textbf{Observation} (Monotonicity of \(\beta_{\mathrm{eff}}\) in \(\rho\)).
Exact enumeration (\(D = 5\)--\(8\)) and Monte Carlo (\(D = 5\)--\(50\),
\(10^5\) samples) establish:

\begin{longtable}[]{@{}llll@{}}
\toprule\noalign{}
\(\rho\) & \(\langle\tilde{r}\rangle\) & \(\beta_{\mathrm{eff}}\) & Regime \\
\midrule\noalign{}
\endhead
\bottomrule\noalign{}
\endlastfoot
0.0 & 0.5006 & 0.64 & Sub-GOE \\
0.3 & 0.5131 & 0.76 & Sub-GOE \\
\textbf{0.5} & \textbf{0.5370} & \textbf{1.01} & \textbf{GOE} \\
0.7 & 0.5861 & 1.70 & GOE--GUE \\
\end{longtable}

At \(\rho = 1/2\) --- the Diaconis--Fulman spectral gap value {[}A; 7{]} --- the
binary model gives \(\beta_{\mathrm{eff}} \approx 1\) (GOE), and the effective
rank is exactly \(g(1/2) = 3/5\) by Proposition 2. The coincidence of the GOE
crossing point with the carry chain's spectral gap is the central result linking
the abstract Markov theory to the observed universality class.

\subsubsection{4.5 Real Eigenvalue Enhancement (C13, Prop.
3)}\label{real-eigenvalue-enhancement-c13-prop.-3}

The fraction of real eigenvalues provides the most direct signature of effective
time-reversal symmetry (TRS): real eigenvalues correspond to eigenvectors that
can be chosen real, the hallmark of the orthogonal class.

\textbf{Observation 3} (Real eigenvalue enhancement). For \(D\)-dimensional
companion matrices with binary Markov entries and correlation \(\rho\)
(established computationally; exact enumeration for \(D = 5\)--\(8\), Monte
Carlo for \(D = 5\)--\(50\)):

\begin{enumerate}
\def\labelenumi{(\alph{enumi})}
\item
  The real eigenvalue fraction \(f_{\mathrm{real}}(\rho)\) is monotonically
  increasing in \(\rho\).
\item
  At \(\rho = 1/2\), \(f_{\mathrm{real}}\) is enhanced by \({\approx}40\%\) over
  the i.i.d. baseline (\(D = 20\)). For actual carry matrices (20-bit primes),
  the enhancement is \({\approx}195\%\) (\(Z = 49.1\)).
\item
  \(f_{\mathrm{real}}\) is strongly anti-correlated with the effective rank: .
\end{enumerate}

\[\mathrm{Corr}(f_{\mathrm{real}}, g(\rho)) = -0.93\]

The anti-correlation (c) confirms the rank-reduction mechanism: as correlation
\(\rho\) increases, the effective rank \(g(\rho)\) decreases (Proposition 2),
and the real eigenvalue fraction increases. Fewer independent perturbation
directions make the companion matrix ``more real,'' consistent with the GOE
transition. The enhancement is stronger for actual carry matrices than for the
binary model because the carry entry distribution is richer than \(\{0,1\}\),
amplifying the correlation effect.

The ratio

\[f_{\mathrm{real}}(\rho{=}0.5) / f_{\mathrm{real}}(\rho{=}0) \approx 1.4\]

is stable across dimensions \(D = 8\) to \(D = 50\). This connects to the real
Ginibre ensemble literature (Edelman {[}13{]}; Akemann, Phillips, and Sommers
{[}14{]}): for real random matrices, the expected number of real eigenvalues
scales as \(\sqrt{2D/\pi}\), and the Markov correlation enhances this by a
factor that depends on \(g(\rho)\).

\subsubsection{4.6 Number Variance Confirmation
(C10)}\label{number-variance-confirmation-c10}

The number variance \(\Sigma^2(L)\) provides independent confirmation:

\begin{longtable}[]{@{}
  >{\raggedright\arraybackslash}p{(\columnwidth - 8\tabcolsep) * \real{0.1190}}
  >{\raggedright\arraybackslash}p{(\columnwidth - 8\tabcolsep) * \real{0.1429}}
  >{\raggedright\arraybackslash}p{(\columnwidth - 8\tabcolsep) * \real{0.1429}}
  >{\raggedright\arraybackslash}p{(\columnwidth - 8\tabcolsep) * \real{0.2619}}
  >{\raggedright\arraybackslash}p{(\columnwidth - 8\tabcolsep) * \real{0.3333}}@{}}
\toprule\noalign{}
\begin{minipage}[b]{\linewidth}\raggedright
\(L\)
\end{minipage} & \begin{minipage}[b]{\linewidth}\raggedright
\(\Sigma^2_{\text{Markov}}\)
\end{minipage} & \begin{minipage}[b]{\linewidth}\raggedright
\(\Sigma^2_{\text{i.i.d.}}\)
\end{minipage} & \begin{minipage}[b]{\linewidth}\raggedright
\(Z\)-score
\end{minipage} & \begin{minipage}[b]{\linewidth}\raggedright
GOE direction?
\end{minipage} \\
\midrule\noalign{}
\endhead
\bottomrule\noalign{}
\endlastfoot
0.5 & 0.334 & 0.317 & +6.3 & Yes
(\(\Sigma^2_{\text{GOE}} > \Sigma^2_{\text{GUE}}\)) \\
1.0 & 0.566 & 0.484 & +24.5 & Yes \\
2.0 & 0.357 & 0.432 & \(-21.8\) & No (reversal) \\
\end{longtable}

At short range (\(L \leq 1\)):
\(\Sigma^2_{\text{Markov}} > \Sigma^2_{\text{i.i.d.}}\) with high significance,
confirming GOE. The carry companion matrices exhibit local level repulsion
consistent with the Wigner surmise for GOE at short ranges, while deviating at
long ranges (\(L \geq 2\)) due to the sparse companion structure. This is
expected: the companion matrix \(M = J + v \cdot e_D^T\) is a rank-1
perturbation of the shift, and universality results for such structures are not
established.

\begin{center}\rule{0.5\linewidth}{0.5pt}\end{center}

\newpage

\subsection{5. Product Zeros: GUE-like}\label{product-zeros-gue-like}

\begin{longtable}[]{@{}
  >{\raggedright\arraybackslash}p{(\columnwidth - 6\tabcolsep) * \real{0.1892}}
  >{\raggedright\arraybackslash}p{(\columnwidth - 6\tabcolsep) * \real{0.3243}}
  >{\raggedright\arraybackslash}p{(\columnwidth - 6\tabcolsep) * \real{0.3243}}
  >{\raggedright\arraybackslash}p{(\columnwidth - 6\tabcolsep) * \real{0.1622}}@{}}
\toprule\noalign{}
\begin{minipage}[b]{\linewidth}\raggedright
Level
\end{minipage} & \begin{minipage}[b]{\linewidth}\raggedright
\(L^2\)(GOE)
\end{minipage} & \begin{minipage}[b]{\linewidth}\raggedright
\(L^2\)(GUE)
\end{minipage} & \begin{minipage}[b]{\linewidth}\raggedright
Best
\end{minipage} \\
\midrule\noalign{}
\endhead
\bottomrule\noalign{}
\endlastfoot
Individual \(M_l\) & 0.331 & 0.409 & GOE \\
Product \(\prod_l\) & 0.512 & 0.494 & weakly GUE-like (3.5\% difference, not
statistically decisive) \\
Known \(\zeta\) zeros & 0.255 & 0.200 & \textbf{GUE} \\
\end{longtable}

The complex weight \(l^{-it}\) in the Euler product introduces incommensurable
phases that break the real-symmetry structure. This is \textbf{analogous} to the
Berry--Keating mechanism {[}10{]}, but with a fundamental distinction:
Berry--Keating postulate a single Hermitian operator; we have an ensemble of
finite-dimensional real matrices whose product is a scalar function. The
GUE-like statistics of product zeros could arise from a central-limit-theorem
effect applied to phase-rotated contributions.

\begin{center}\rule{0.5\linewidth}{0.5pt}\end{center}

\newpage

\subsection{6. Discussion}\label{discussion}

\begin{longtable}[]{@{}
  >{\raggedright\arraybackslash}p{(\columnwidth - 4\tabcolsep) * \real{0.2414}}
  >{\raggedright\arraybackslash}p{(\columnwidth - 4\tabcolsep) * \real{0.3793}}
  >{\raggedright\arraybackslash}p{(\columnwidth - 4\tabcolsep) * \real{0.3793}}@{}}
\toprule\noalign{}
\begin{minipage}[b]{\linewidth}\raggedright
Level
\end{minipage} & \begin{minipage}[b]{\linewidth}\raggedright
Statistics
\end{minipage} & \begin{minipage}[b]{\linewidth}\raggedright
Mechanism
\end{minipage} \\
\midrule\noalign{}
\endhead
\bottomrule\noalign{}
\endlastfoot
Individual \(M_l\) & GOE-like & Markov correlation → rank reduction (Prop. 2);
positive autocorrelation (Prop. 1); real eigenvalue enhancement (Prop. 3) \\
i.i.d. entries & GUE-like & Decorrelation (§3.1); full-rank perturbation \\
Transition & Smooth, \(\beta_{\mathrm{eff}}(\lambda)\) & Continuous
interpolation (§4.1); \(\Sigma^2\) confirms at short range (§4.6) \\
Product \(\prod_l\) & GUE-like (marginal) & Complex phases from \(l^{-it}\) \\
\end{longtable}

\subsubsection{6.1 Connection to Boundary Spectral Structure
{[}E{]}}\label{connection-to-boundary-spectral-structure-e}

The overall convergence rate of the trace anomaly is governed by the
Diaconis--Fulman eigenvalue \(\lambda_2 = 1/b\) {[}A; E; G, §4.5{]}. This same
eigenvalue determines the effective Markov correlation in the binary model
(§4.4), establishing a direct link between the spectral theory of the carry
operator and the observed universality class of the carry companion matrices.
The anti-correlation conjecture (\(\alpha_k \to (b-1)/b\) for all \(k \geq 2\);
{[}A, Conjecture 4{]}) would, if proved, constrain the asymptotic Markov
correlation and hence the limiting \(\beta_{\mathrm{eff}}\). No rigorous
analytical connection between the spectral gap and the Dyson index \(\beta\) has
been established; this remains an open problem.

\subsubsection{Open Questions}\label{open-questions}

\begin{enumerate}
\def\labelenumi{\arabic{enumi}.}
\item
  \textbf{Sub-GOE effect:} The spacing ratio at \(D = 50\) falls below the GOE
  reference (\(\langle\tilde{r}\rangle = 0.495\) vs \(0.531\)). Does this
  indicate \(\beta < 1\) for sparse companion matrices, or is it a finite-size
  correction? The sparse companion structure (\(D-1\) deterministic entries, 1
  random column) may produce statistics outside standard \(\beta\)-ensembles.
\item
  \textbf{Analytical bound:} The computational bound (§4.3) establishes
\end{enumerate}

\[\beta_{\mathrm{eff}} \leq 2 - \varepsilon\]

numerically. An analytical proof --- even for a weak bound like
\(\varepsilon > 0\) when Markov lag-1 autocorrelation is positive --- would
require relating the effective rank \(g(\rho)\) (Proposition 2) to the level
repulsion exponent \(\beta\), an original contribution to the theory of sparse
non-Hermitian ensembles.

\begin{enumerate}
\def\labelenumi{\arabic{enumi}.}
\setcounter{enumi}{2}
\item
  \textbf{Product convergence:} How many Euler factors are needed for the
  GOE→GUE transition in the product spectrum?
\item
  \textbf{Quantitative real eigenvalue prediction:} Observation 3 establishes
  that Markov correlation increases \(f_{\mathrm{real}}\) by \({\sim}40\%\)
  (binary model) to \({\sim}195\%\) (actual carries). A closed-form prediction
  of the enhancement factor as a function of \(g(\rho)\) --- connecting to the
  real Ginibre ensemble literature (Edelman {[}13{]}; Akemann, Phillips, and
  Sommers {[}14{]}) --- would close the analytical loop between the effective
  rank mechanism and the TRS signature.
\end{enumerate}

\begin{center}\rule{0.5\linewidth}{0.5pt}\end{center}

\newpage

\subsection{7. Reproducibility}\label{reproducibility}

13 experiment scripts in \texttt{experiments/}:
\texttt{C01\_goe\_gue\_unitary\_transition.py},
\texttt{C02\_goe\_gue\_finite\_size.py},
\texttt{C03\_factorials\_goe\_gue\_transition.py},
\texttt{C04\_gue\_correlation.py}, \texttt{C05\_goe\_gue\_scaling\_limit.py},
\texttt{C06\_analytical\_ensemble\_structure.py},
\texttt{C07\_goe\_spacing\_ratio.py}, \texttt{C08\_beta\_interpolation.py},
\texttt{C09\_symmetry\_mechanism.py}, \texttt{C10\_number\_variance.py},
\texttt{C11\_beta\_bound.py}, \texttt{C12\_analytical\_beta\_lemma.py},
\texttt{C13\_analytical\_foundations.py} (Props. 1--3). Shared utilities in
\texttt{src/carry\_utils.py}.

Requirements: Python 3.8+, NumPy, SciPy.

\begin{center}\rule{0.5\linewidth}{0.5pt}\end{center}

\newpage

\subsection{8. References}\label{references}

\begin{enumerate}
\def\labelenumi{\arabic{enumi}.}
\tightlist
\item
  H. L. Montgomery, ``The pair correlation of zeros of the zeta function,''
  \emph{Proc. Symp. Pure Math.} 24, 181--193, 1973.
\item
  N. M. Katz and P. Sarnak, \emph{Random Matrices, Frobenius Eigenvalues, and
  Monodromy}, AMS Colloquium Publications, vol.~45, 1999.
\item
  L. Erdős, B. Schlein, and H.-T. Yau, ``Universality of random matrices and
  local relaxation flow,'' \emph{Invent. Math.} 185(1), 75--119, 2011.
\item
  L. Erdős and H.-T. Yau, ``Universality of local spectral statistics of random
  matrices,'' \emph{Bull. Amer. Math. Soc.} 49(3), 377--414, 2012.
\item
  T. Tao and V. Vu, ``Random matrices: Universality of local spectral statistics
  of non-Hermitian matrices,'' \emph{Ann. Probab.} 43(2), 782--874, 2015.
\item
  I. Dumitriu and A. Edelman, ``Matrix models for beta ensembles,'' \emph{J.
  Math. Phys.} 43(11), 5830--5847, 2002.
\item
  P. Diaconis and J. Fulman, ``Carries, shuffling, and symmetric functions,''
  \emph{Adv. Appl. Math.} 43(2), 176--196, 2009.
\item
  U. Grenander and G. Szegő, \emph{Toeplitz Forms and Their Applications},
  University of California Press, 1958 (repr. Chelsea, 1984).
\item
  Y. Y. Atas, E. Bogomolny, O. Giraud, and G. Roux, ``Distribution of the ratio
  of consecutive level spacings in random matrix ensembles,'' \emph{Phys.
  Rev.~Lett.} 110, 084101, 2013.
\item
  M. V. Berry and J. P. Keating, ``The Riemann zeros and eigenvalue
  asymptotics,'' \emph{SIAM Review} 41(2), 236--266, 1999.
\item
  M. L. Mehta, \emph{Random Matrices}, 3rd ed., Academic Press, 2004.
\item
  F. J. Dyson, ``Statistical theory of the energy levels of complex systems,''
  \emph{J. Math. Phys.} 3, 140--175, 1962.
\item
  A. Edelman, ``The probability that a random real Gaussian matrix has \(k\)
  real eigenvalues, related distributions, and the circular law,'' \emph{J.
  Multivar. Anal.} 60(2), 203--232, 1997.
\item
  G. Akemann, M. J. Phillips, and H.-J. Sommers, ``The chiral Gaussian
  two-matrix ensemble of real asymmetric matrices,'' \emph{J. Phys. A: Math.
  Theor.} 43, 085211, 2010.
\item
  {[}A{]} Companion paper: ``Spectral Theory of Carries in Positional
  Multiplication,'' this series.
\item
  {[}E{]} Companion paper: ``The Trace Anomaly of Binary Multiplication,'' this
  series.
\item
  {[}F{]} Companion paper: ``Exact Covariance Structure of Binary Carry
  Chains,'' this series.
\item
  {[}G{]} Companion paper: ``The Angular Uniqueness of Base 2 in Positional
  Multiplication,'' this series.
\item
  {[}P1{]} Companion paper: ``π from Pure Arithmetic: A Spectral Phase
  Transition in the Binary Carry Bridge,'' this series.
\end{enumerate}

\begin{center}\rule{0.5\linewidth}{0.5pt}\end{center}

\emph{CC BY 4.0. Code: MIT License.}

\end{document}
